% ═══════════════════════════════════════════════════════════════
% MUSTERLÖSUNG ML-08b: El tiempo (Das Wetter)
% ═══════════════════════════════════════════════════════════════
% Fach: Spanisch | Klasse: 7 | Sequenz: 08 (Vacaciones)
% Fachfarbe: #c41e3a (Karminrot)
% Lösungsfarbe: ROT (deutlich sichtbar)
%
% WICHTIG: Identisches Layout wie AB, nur mit roten Lösungen
% ═══════════════════════════════════════════════════════════════

\documentclass[11pt, a4paper]{article}

% ═══════════════════════════════════════════════════════════════
% SW-MODUS TOGGLE
% ═══════════════════════════════════════════════════════════════
\newif\ifswmodus
\swmodusfalse    % Standard: Farbe

% ═══════════════════════════════════════════════════════════════
% PAKETE
% ═══════════════════════════════════════════════════════════════

\usepackage[utf8]{inputenc}
\usepackage[T1]{fontenc}
\usepackage[ngerman]{babel}
\usepackage{helvet}
\renewcommand{\familydefault}{\sfdefault}

\usepackage[margin=2cm]{geometry}
\usepackage{enumitem}
\usepackage{tabularx}
\usepackage{booktabs}
\usepackage{multirow}
\usepackage{graphicx}

\usepackage{xcolor}
\usepackage{framed}
\usepackage{tcolorbox}
\tcbuselibrary{skins}

\usepackage{fancyhdr}
\usepackage{lastpage}
\usepackage{needspace}

% ═══════════════════════════════════════════════════════════════
% FARBEN
% ═══════════════════════════════════════════════════════════════

\definecolor{spanisch-color}{HTML}{c41e3a}
\definecolor{grau-color}{HTML}{888888}
\definecolor{hellgrau-color}{HTML}{f5f5f5}
\definecolor{orange-color}{HTML}{b35c00}
\definecolor{loesung-color}{HTML}{c62828}

\definecolor{sw-dark}{HTML}{000000}
\definecolor{sw-gray}{HTML}{666666}
\definecolor{sw-white}{HTML}{ffffff}

\ifswmodus
    \colorlet{spanisch}{sw-dark}
    \colorlet{grau}{sw-gray}
    \colorlet{hellgrau}{sw-white}
    \colorlet{orange}{sw-dark}
    \colorlet{loesung}{sw-dark}
\else
    \colorlet{spanisch}{spanisch-color}
    \colorlet{grau}{grau-color}
    \colorlet{hellgrau}{hellgrau-color}
    \colorlet{orange}{orange-color}
    \colorlet{loesung}{loesung-color}
\fi

\newcommand{\fachfarbe}{spanisch}

% ═══════════════════════════════════════════════════════════════
% VARIABLEN
% ═══════════════════════════════════════════════════════════════

\newcommand{\thema}{El tiempo -- Das Wetter}
\newcommand{\sequenz}{08}
\newcommand{\block}{b}
\newcommand{\fach}{Spanisch}
\newcommand{\klasse}{7}
\newcommand{\maxpunkte}{20}
\newcommand{\datum}{\today}

% ═══════════════════════════════════════════════════════════════
% KOPF- UND FUSSZEILE
% ═══════════════════════════════════════════════════════════════

\pagestyle{fancy}
\fancyhf{}
\renewcommand{\headrulewidth}{0.4pt}
\renewcommand{\footrulewidth}{0.4pt}

\lhead{\fach{} -- Klasse \klasse}
\chead{\textcolor{loesung}{\textbf{ML-\sequenz\block{} (MUSTERLÖSUNG)}}}
\rhead{\datum}

\lfoot{}
\cfoot{}
\rfoot{Seite \thepage\ von \pageref{LastPage}}

% ═══════════════════════════════════════════════════════════════
% BEFEHLE
% ═══════════════════════════════════════════════════════════════

\newcommand{\punktebox}[1]{\hfill\fbox{\small \_/#1}}

\newcommand{\loesung}[1]{\textcolor{loesung}{\textbf{#1}}}

\newcommand{\luecke}[1]{\rule{#1}{0.4pt}}

\newtcolorbox{merkkasten}[1][]{
    colback=hellgrau,
    colframe=\fachfarbe,
    colbacktitle=white,
    coltitle=\fachfarbe,
    boxrule=1pt,
    arc=0pt,
    left=8pt, right=8pt, top=8pt, bottom=8pt,
    fonttitle=\bfseries,
    attach boxed title to top left={yshift=-2mm, xshift=5mm},
    boxed title style={boxrule=0pt, colback=white},
    title=#1
}

\newtcolorbox{vokabelkasten}{
    colback=white,
    colframe=\fachfarbe,
    boxrule=0.5pt,
    arc=0pt,
    left=5pt, right=5pt, top=5pt, bottom=5pt
}

\newtcolorbox{hinweis}{
    colback=white,
    colframe=orange,
    colbacktitle=white,
    coltitle=orange,
    boxrule=1pt,
    arc=0pt,
    left=5pt, right=5pt, top=5pt, bottom=5pt,
    fonttitle=\bfseries,
    attach boxed title to top left={yshift=-2mm, xshift=5mm},
    boxed title style={boxrule=0pt, colback=white},
    title=Hinweis
}

\newenvironment{aufgabe}[1]{%
    \needspace{5\baselineskip}%
    \nopagebreak[4]%
    \vspace{10pt}
    \noindent\textbf{\textcolor{\fachfarbe}{Aufgabe #1}}
    \nopagebreak[4]%
    \vspace{5pt}
    \nopagebreak[4]%
    \begin{list}{}{\leftmargin=0pt}
    \item[]
}{%
    \end{list}
}

\newcommand{\es}[1]{\textcolor{\fachfarbe}{\textbf{#1}}}

% ═══════════════════════════════════════════════════════════════
% DOKUMENT
% ═══════════════════════════════════════════════════════════════

\begin{document}

% ═══════════════════════════════════════════════════════════════
% KOPFBEREICH
% ═══════════════════════════════════════════════════════════════

\begin{center}
    {\Large\bfseries\textcolor{\fachfarbe}{Arbeitsblatt: \thema}}\\[0.3em]
    {\large\textcolor{loesung}{\textbf{--- MUSTERLÖSUNG ---}}}
\end{center}

\vspace{5pt}

\noindent
\begin{tabularx}{\textwidth}{|X|X|X|}
\hline
\textbf{Name:} \textcolor{loesung}{Musterschüler/in} &
\textbf{Klasse:} \klasse &
\textbf{Datum:} \datum \\
\hline
\end{tabularx}

\vspace{15pt}

% ═══════════════════════════════════════════════════════════════
% MERKKASTEN 1: Das Wetter (mit Lösungen)
% ═══════════════════════════════════════════════════════════════

\begin{merkkasten}[Merkkasten 1: El tiempo -- Das Wetter]
\textbf{Wetter mit \es{hacer} (machen):}

\vspace{0.5em}

\begin{tabular}{ll}
\es{Hace sol.} & = \loesung{Die Sonne scheint.} \\[0.8em]
\es{Hace calor.} & = \loesung{Es ist heiß/warm.} \\[0.8em]
\es{Hace frío.} & = \loesung{Es ist kalt.} \\[0.8em]
\es{Hace viento.} & = \loesung{Es ist windig.} \\[0.8em]
\es{Hace buen tiempo.} & = \loesung{Das Wetter ist gut.} \\[0.8em]
\es{Hace mal tiempo.} & = \loesung{Das Wetter ist schlecht.} \\
\end{tabular}

\vspace{1em}

\textbf{Wetter mit Verben und \es{estar}:}

\vspace{0.5em}

\begin{tabular}{ll}
\es{Llueve.} & = \loesung{Es regnet.} \\[0.8em]
\es{Nieva.} & = \loesung{Es schneit.} \\[0.8em]
\es{Está nublado.} & = \loesung{Es ist bewölkt.} \\
\end{tabular}

\end{merkkasten}

\vspace{10pt}

% ═══════════════════════════════════════════════════════════════
% MERKKASTEN 2: Die Jahreszeiten (mit Lösungen)
% ═══════════════════════════════════════════════════════════════

\begin{merkkasten}[Merkkasten 2: Las estaciones -- Die Jahreszeiten]

\vspace{0.5em}

\begin{tabular}{ll}
\es{la primavera} & = \loesung{der Frühling} \\[0.8em]
\es{el verano} & = \loesung{der Sommer} \\[0.8em]
\es{el otoño} & = \loesung{der Herbst} \\[0.8em]
\es{el invierno} & = \loesung{der Winter} \\
\end{tabular}

\vspace{0.5em}

\textbf{Fragen:} \es{¿Qué tiempo hace?} = \loesung{Wie ist das Wetter?}
\end{merkkasten}

\vspace{15pt}

% ═══════════════════════════════════════════════════════════════
% AUFGABE 1: Zuordnung (mit Lösungen)
% ═══════════════════════════════════════════════════════════════

\begin{aufgabe}{1}
Verbinde Spanisch und Deutsch. Ziehe Linien. \punktebox{4}
\end{aufgabe}

\begin{center}
\begin{tabular}{rcl}
\es{Hace sol.} & \loesung{---} & Es ist windig. \loesung{(= Hace viento)} \\[0.8em]
\es{Hace frío.} & \loesung{---} & Es regnet. \loesung{(= Llueve)} \\[0.8em]
\es{Hace viento.} & \loesung{---} & Die Sonne scheint. \loesung{(= Hace sol)} \\[0.8em]
\es{Llueve.} & \loesung{---} & Es ist kalt. \loesung{(= Hace frío)} \\
\end{tabular}
\end{center}

\loesung{Lösung: Hace sol $\rightarrow$ Die Sonne scheint | Hace frío $\rightarrow$ Es ist kalt | Hace viento $\rightarrow$ Es ist windig | Llueve $\rightarrow$ Es regnet}

\vspace{15pt}

% ═══════════════════════════════════════════════════════════════
% AUFGABE 2: Lückentext (mit Lösungen)
% ═══════════════════════════════════════════════════════════════

\begin{aufgabe}{2}
Ergänze die Lücken mit dem passenden Wetterwort. \punktebox{4}
\end{aufgabe}

\begin{vokabelkasten}
\textit{Wortliste:} calor | nublado | nieva | primavera
\end{vokabelkasten}

\vspace{0.5em}

\noindent
a) En \es{verano} hace mucho \loesung{calor}. \\[0.8em]
b) En \es{invierno} a veces \loesung{nieva} en las montañas. \\[0.8em]
c) Hoy el cielo está \loesung{nublado}, no hay sol. \\[0.8em]
d) En \loesung{primavera} hace buen tiempo y hay muchas flores. \\

\vspace{15pt}

% ═══════════════════════════════════════════════════════════════
% AUFGABE 3: Jahreszeiten und Wetter (mit Lösungen)
% ═══════════════════════════════════════════════════════════════

\begin{aufgabe}{3}
Welches Wetter passt zu welcher Jahreszeit? Schreibe Sätze. \punktebox{4}
\end{aufgabe}

\noindent
\textbf{Beispiel:} \textit{En verano hace calor y hace sol.}

\vspace{0.5em}

\noindent
a) \textbf{En invierno:}

\loesung{En invierno hace frío y a veces nieva.}

\vspace{0.8em}

b) \textbf{En otoño:}

\loesung{En otoño hace viento y llueve.}

\vspace{0.8em}

c) \textbf{En primavera:}

\loesung{En primavera hace buen tiempo y hace sol.}

\textit{\textcolor{grau}{Bewertung: Je 1P für korrekten Wetterausdruck + 1P für vollständigen Satz}}

\vspace{15pt}

% ═══════════════════════════════════════════════════════════════
% AUFGABE 4: Mini-Wetterbericht (Beispiellösung)
% ═══════════════════════════════════════════════════════════════

\begin{aufgabe}{4}
Du bist im Urlaub in Barcelona. Schreibe einen kurzen Wetterbericht (mind. 4 Sätze). \punktebox{8}
\end{aufgabe}

\begin{hinweis}
Verwende: Jahreszeit, Wetter, was du machst. \\
\textbf{Hilfe:} Estoy en... | Es... (Jahreszeit) | Hace... | Hoy... | Voy a...
\end{hinweis}

\vspace{0.5em}

\textbf{Beispiellösung:}

\loesung{¡Hola! Estoy en Barcelona.}

\loesung{Es verano y hace mucho calor.}

\loesung{Hoy hace sol y hace buen tiempo.}

\loesung{Voy a la playa porque hace calor.}

\loesung{¡Me gusta el verano en España!}

\vspace{0.5em}
\textit{\textcolor{grau}{Bewertung: Ort (1P), Jahreszeit (1P), 2 Wetterausdrücke (2P), Aktivität (2P), sprachliche Korrektheit (2P)}}

% ═══════════════════════════════════════════════════════════════
% PUNKTEÜBERSICHT
% ═══════════════════════════════════════════════════════════════

\vspace{20pt}
\hrule
\vspace{10pt}

\begin{flushright}
\textbf{Punkte:} \loesung{\maxpunkte} / \maxpunkte
\end{flushright}

\end{document}
